% Options for packages loaded elsewhere
\PassOptionsToPackage{unicode}{hyperref}
\PassOptionsToPackage{hyphens}{url}
\PassOptionsToPackage{dvipsnames,svgnames,x11names}{xcolor}
%
\documentclass[
  11pt,
]{article}
\usepackage{amsmath,amssymb}
\usepackage{iftex}
\ifPDFTeX
  \usepackage[T1]{fontenc}
  \usepackage[utf8]{inputenc}
  \usepackage{textcomp} % provide euro and other symbols
\else % if luatex or xetex
  \usepackage{unicode-math} % this also loads fontspec
  \defaultfontfeatures{Scale=MatchLowercase}
  \defaultfontfeatures[\rmfamily]{Ligatures=TeX,Scale=1}
\fi
\usepackage{lmodern}
\ifPDFTeX\else
  % xetex/luatex font selection
\fi
% Use upquote if available, for straight quotes in verbatim environments
\IfFileExists{upquote.sty}{\usepackage{upquote}}{}
\IfFileExists{microtype.sty}{% use microtype if available
  \usepackage[]{microtype}
  \UseMicrotypeSet[protrusion]{basicmath} % disable protrusion for tt fonts
}{}
\makeatletter
\@ifundefined{KOMAClassName}{% if non-KOMA class
  \IfFileExists{parskip.sty}{%
    \usepackage{parskip}
  }{% else
    \setlength{\parindent}{0pt}
    \setlength{\parskip}{6pt plus 2pt minus 1pt}}
}{% if KOMA class
  \KOMAoptions{parskip=half}}
\makeatother
\usepackage{xcolor}
\usepackage[margin=2.2cm]{geometry}
\setlength{\emergencystretch}{3em} % prevent overfull lines
\providecommand{\tightlist}{%
  \setlength{\itemsep}{0pt}\setlength{\parskip}{0pt}}
\setcounter{secnumdepth}{5}
\usepackage{lmodern}
\usepackage[T1]{fontenc}
\usepackage{amsmath,amssymb}
\usepackage{newunicodechar}
\newunicodechar{}{\ensuremath{\bar\nu}}
\newunicodechar{}{\ensuremath{\hat\nu}}
\newunicodechar{}{\ensuremath{\tilde\sigma}}
\newunicodechar{}{\ensuremath{\bar v}}
\newunicodechar{}{\ensuremath{\hat\sigma}}
\newunicodechar{}{\ensuremath{\bar X}}
\newunicodechar{}{\ensuremath{\hat X}}
\newunicodechar{}{\ensuremath{\tilde W}}
\newunicodechar{}{\ensuremath{\dot\alpha}}
\newunicodechar{}{\ensuremath{\dot\beta}}
\newunicodechar{}{\ensuremath{\hat f}}
\newunicodechar{}{\ensuremath{\bar t}}
\newunicodechar{}{\ensuremath{\hat V}}
\newunicodechar{Ĉ}{\ensuremath{\hat C}}
\newunicodechar{ĉ}{\ensuremath{\hat c}}
\newunicodechar{ġ}{\ensuremath{\dot g}}
\newunicodechar{ᵀ}{\ensuremath{{}^{\mathsf{T}}}}
\newunicodechar{ℝ}{\ensuremath{\mathbb{R}}}
\newunicodechar{−}{\ensuremath{-}}
\newunicodechar{²}{\ensuremath{{}^{2}}}
\newunicodechar{³}{\ensuremath{{}^{3}}}
\newunicodechar{¹}{\ensuremath{{}^{1}}}
\newunicodechar{⁴}{\ensuremath{{}^{4}}}
\newunicodechar{⁸}{\ensuremath{{}^{8}}}
\newunicodechar{⁺}{\ensuremath{{}^{+}}}
\newunicodechar{⁻}{\ensuremath{{}^{-}}}
\newunicodechar{₂}{\ensuremath{{}_{2}}}
\newunicodechar{½}{\ensuremath{\tfrac12}}
\newunicodechar{¼}{\ensuremath{\tfrac14}}
\newunicodechar{·}{\ensuremath{\cdot}}
\newunicodechar{×}{\ensuremath{\times}}
\newunicodechar{÷}{\ensuremath{\div}}
\newunicodechar{±}{\ensuremath{\pm}}
\newunicodechar{σ}{\ensuremath{\sigma}}
\newunicodechar{φ}{\ensuremath{\varphi}}
\newunicodechar{ν}{\ensuremath{\nu}}
\newunicodechar{τ}{\ensuremath{\tau}}
\newunicodechar{π}{\ensuremath{\pi}}
\newunicodechar{ρ}{\ensuremath{\rho}}
\newunicodechar{λ}{\ensuremath{\lambda}}
\newunicodechar{μ}{\ensuremath{\mu}}
\newunicodechar{κ}{\ensuremath{\kappa}}
\newunicodechar{ω}{\ensuremath{\omega}}
\newunicodechar{β}{\ensuremath{\beta}}
\newunicodechar{α}{\ensuremath{\alpha}}
\newunicodechar{γ}{\ensuremath{\gamma}}
\newunicodechar{θ}{\ensuremath{\theta}}
\newunicodechar{η}{\ensuremath{\eta}}
\newunicodechar{χ}{\ensuremath{\chi}}
\newunicodechar{ψ}{\ensuremath{\psi}}
\newunicodechar{ξ}{\ensuremath{\xi}}
\newunicodechar{δ}{\ensuremath{\delta}}
\newunicodechar{ε}{\ensuremath{\varepsilon}}
\newunicodechar{Σ}{\ensuremath{\Sigma}}
\newunicodechar{Δ}{\ensuremath{\Delta}}
\newunicodechar{Λ}{\ensuremath{\Lambda}}
\newunicodechar{Π}{\ensuremath{\Pi}}
\newunicodechar{Ξ}{\ensuremath{\Xi}}
\newunicodechar{Γ}{\ensuremath{\Gamma}}
\newunicodechar{Φ}{\ensuremath{\Phi}}
\newunicodechar{Θ}{\ensuremath{\Theta}}
\newunicodechar{Ψ}{\ensuremath{\Psi}}
\newunicodechar{∫}{\ensuremath{\int}}
\newunicodechar{∂}{\ensuremath{\partial}}
\newunicodechar{√}{\ensuremath{\surd}}
\newunicodechar{∞}{\ensuremath{\infty}}
\newunicodechar{≈}{\ensuremath{\approx}}
\newunicodechar{≥}{\ensuremath{\ge}}
\newunicodechar{≤}{\ensuremath{\le}}
\newunicodechar{≠}{\ensuremath{\ne}}
\newunicodechar{≡}{\ensuremath{\equiv}}
\newunicodechar{→}{\ensuremath{\to}}
\newunicodechar{←}{\ensuremath{\leftarrow}}
\newunicodechar{⇒}{\ensuremath{\Rightarrow}}
\newunicodechar{⇔}{\ensuremath{\Leftrightarrow}}
\newunicodechar{↦}{\ensuremath{\mapsto}}
\newunicodechar{∈}{\ensuremath{\in}}
\newunicodechar{∝}{\ensuremath{\propto}}
\newunicodechar{∏}{\ensuremath{\prod}}
\newunicodechar{⊥}{\ensuremath{\perp}}
\newunicodechar{‖}{\ensuremath{\|}}
\newunicodechar{⟨}{\ensuremath{\langle}}
\newunicodechar{⟩}{\ensuremath{\rangle}}
\newunicodechar{∩}{\ensuremath{\cap}}
\newunicodechar{⊆}{\ensuremath{\subseteq}}
\newunicodechar{…}{\ensuremath{\ldots}}
\newunicodechar{̄}{}
\newunicodechar{̂}{}
\newunicodechar{̃}{}
\newunicodechar{̇}{}
\ifLuaTeX
  \usepackage{selnolig}  % disable illegal ligatures
\fi
\IfFileExists{bookmark.sty}{\usepackage{bookmark}}{\usepackage{hyperref}}
\IfFileExists{xurl.sty}{\usepackage{xurl}}{} % add URL line breaks if available
\urlstyle{same}
\hypersetup{
  pdftitle={Assignment-Defense Questions},
  pdfauthor={Computational Finance (WI4151), TU Delft --- Prof.~Fang Fang},
  colorlinks=true,
  linkcolor={blue},
  filecolor={Maroon},
  citecolor={Blue},
  urlcolor={Blue},
  pdfcreator={LaTeX via pandoc}}

\title{Assignment-Defense Questions}
\author{Computational Finance (WI4151), TU Delft --- Prof.~Fang Fang}
\date{}

\begin{document}
\maketitle

{
\hypersetup{linkcolor=blue}
\setcounter{tocdepth}{3}
\tableofcontents
}
Code-level questions for each assignment. Expect ``open your file and
explain line X.'' Open the actual file when answering and give the line
number back.

\hypertarget{exercise-set-1}{%
\section{Exercise set 1}\label{exercise-set-1}}

Files: -
\texttt{assignments/exercise-set-1/exercise\_1\_computational\_finance\ (2).py}
-
\texttt{assignments/exercise-set-1/exercise\_4\_computational\_finance\ (2).py}

\hypertarget{exercise-1-implied-volatility}{%
\subsection{Exercise 1 --- implied
volatility}\label{exercise-1-implied-volatility}}

\begin{enumerate}
\def\labelenumi{\arabic{enumi}.}
\item
  Open \texttt{exercise\_1\_computational\_finance\ (2).py} and explain
  lines 13--14 of \(C_bs\): what does the guard
  \texttt{if\ sigma\ \textless{}=\ 0\ or\ T\ \textless{}=\ 0} return,
  and why is \(max(S0 − K·e^{−rT}, 0)\) the right fallback value rather
  than \(0\)? \textbf{A.} The guard returns the option's intrinsic value
  in forward-discounted terms, \(max(S0 − K e^{−rT}, 0)\). When
  \(sigma → 0\) or \(T → 0\) the Black--Scholes formula divides by
  \(sigma√T → 0\) (the \(d1\), \(d2\) definitions on lines 15--16 are
  undefined), so I short-circuit. The right limiting value is not \(0\):
  as \(σ→0\) the stock grows deterministically to \(S0 e^{rT}\), and the
  discounted payoff
  \texttt{e\^{}\{−rT\}\ max(S0\ e\^{}\{rT\}\ −\ K,\ 0)\ =\ max(S0\ −\ K\ e\^{}\{−rT\},\ 0)}.
  That is the no-arbitrage lower bound of a call, so the function stays
  continuous and never returns a value below the bound. Returning plain
  \(0\) would wrongly price a deep-ITM call at zero.
\item
  In \texttt{implied\_volatility\_bisection} (lines 19--41), what does
  the function return on line 29, and under what condition? Explain in
  terms of the sign of \texttt{f\_low·f\_high}. \textbf{A.} Line 29
  returns \texttt{np.nan} when
  \texttt{f\_low\ *\ f\_high\ \textgreater{}\ 0}, i.e.~when
  \texttt{C\_bs(σ\_low)−C\_mkt} and \texttt{C\_bs(σ\_high)−C\_mkt} have
  the same sign. Bisection needs a sign change to guarantee a root in
  \texttt{{[}σ\_low,\ σ\_high{]}} (intermediate value theorem); if both
  endpoints sit on the same side, the market price lies outside the
  achievable BS range on \([1e-6, 5.0]\) (e.g.~a price violating
  no-arbitrage bounds), so no implied vol exists in the bracket and
  returning NaN flags that path rather than producing a garbage number.
\item
  Line 34 has a compound stopping test
  \texttt{abs(f\_mid)\ \textless{}\ tol\ or\ (sigma\_high\ -\ sigma\_low)\ \textless{}\ tol}.
  Which of the two conditions is in \emph{price} units and which is in
  \emph{vol} units? Why can the loop exit on either? \textbf{A.}
  \(abs(f_mid) < tol\) is in price (option-value) units ---
  \texttt{f\_mid\ =\ C\_bs(σ\_mid)−C\_mkt} is a price residual.
  \texttt{(sigma\_high\ −\ sigma\_low)\ \textless{}\ tol} is in
  volatility units --- the width of the bracket. Either is a valid
  termination: the first says the price is matched to tolerance, the
  second says the vol is pinned to tolerance. Both must hold eventually
  since the bracket halves each step; including both protects against
  the flat-vega regime where the price residual stays large per unit of
  σ but the interval has already collapsed.
\item
  \(Vega\) (lines 43--49) returns \texttt{S0*np.sqrt(T)*norm.pdf(d1)}.
  Why is this exactly \texttt{∂C\_BS/∂σ}, and why does it appear in the
  Newton denominator on line 66? \textbf{A.} Differentiating
  \(C_BS = S0 N(d1) − K e^{−rT} N(d2)\) w.r.t. σ, the explicit σ-terms
  via \(d1\), \(d2\) cancel by the identity
  \(S0 φ(d1) = K e^{−rT} φ(d2)\) (where φ = \texttt{norm.pdf}), leaving
  only \texttt{∂C/∂σ\ =\ S0\ φ(d1)\ ∂d1/∂σ\ −\ {[}same\ cancels{]}}, and
  \texttt{∂d1/∂σ}'s clean part gives \(S0 √T φ(d1)\). So Vega is
  \texttt{∂C\_BS/∂σ}. Newton solves \texttt{f(σ)=C\_BS(σ)−C\_mkt=0}, and
  \(f'(σ)=Vega\), so line 66 \texttt{σ\ -=\ f/f\textquotesingle{}} is
  \texttt{σ\ -=\ diff/vega}.
\item
  Explain line 58:
  \texttt{sigma\_initial\ =\ np.sqrt((1/T)*(2*np.abs(np.log(S0/K)+\ r*T)))}.
  What is this initial guess and what does the report claim it
  guarantees? Where in the lectures does it come from? \textbf{A.} This
  is the Brenner--Subrahmanyam / Manaster--Koehler starting point
  \texttt{σ0\ =\ √(2/T\ ·\ \textbar{}ln(S0/K)\ +\ rT\textbar{})}.
  Manaster--Koehler proved this is the σ at which Vega is maximal (the
  inflection of C\_BS in σ), which puts Newton on the side of the curve
  where it converges \emph{monotonically and globally} to the implied
  vol --- that is the guarantee the report cites. It comes from the
  Lecture 3 root-finding slides on Newton for implied vol.
\item
  On line 66, \texttt{sigma\ -=\ diff/vega}. If \(vega\) were
  \textasciitilde0 (deep OTM, short T), what would happen numerically,
  and which method (bisection vs Newton) is safe against this? Which
  line shows the safe behaviour? \textbf{A.} If \(vega ≈ 0\),
  \texttt{diff/vega} blows up (division by near-zero) and Newton takes a
  huge, possibly diverging step or overflows. Bisection is safe because
  it never divides by a derivative --- it only checks signs and halves
  the interval (lines 31--40); its convergence depends only on the
  bracketing, not on Vega. The safe behaviour is on line 34's interval
  test and the halving on line 31.
\item
  Lines 94--96 build \(IV_bisection\) and \(IV_Newton\). What column is
  passed as \(C_mkt\), and why does that choice (mid vs bid/ask)
  directly feed your Exercise 1.3 discrepancy discussion? \textbf{A.}
  \(C_mkt\) is the \('Midpoint'\) column (line 94--95 zip over
  \(df_calls['Midpoint']\)). Mid = (bid+ask)/2, so the inverted IV is a
  single point estimate sitting inside the bid--ask spread.
  Discrepancies versus the CSV reference IV (and the bid/ask one would
  get from the quoted spread) come from this choice: the bid--ask spread
  maps to an IV band, and using mid collapses it to one number --- that
  band width, plus discretisation of strikes, is exactly the source of
  the 1.3 mismatch.
\item
  Roughly how many bisection iterations are needed to shrink the bracket
  \([1e-6, 5.0]\) below \(tol = 1e-6\)? Why is \(max_iter = 200\) (line
  90) more than enough? \textbf{A.} The bracket has width ≈ 5, and each
  step halves it, so I need
  \texttt{5/2\^{}n\ \textless{}\ 1e-6\ ⇒\ 2\^{}n\ \textgreater{}\ 5e6\ ⇒\ n\ \textgreater{}\ log2(5e6)\ ≈\ 22.3},
  so about 23 iterations. \(max_iter = 200\) is \textasciitilde10× that,
  so the iteration cap is never the binding constraint; the tolerance
  test on line 34 always fires first.
\end{enumerate}

\hypertarget{exercise-4-monte-carlo-basket-call}{%
\subsection{Exercise 4 --- Monte Carlo basket
call}\label{exercise-4-monte-carlo-basket-call}}

\begin{enumerate}
\def\labelenumi{\arabic{enumi}.}
\setcounter{enumi}{8}
\item
  Open \texttt{exercise\_4\_computational\_finance\ (2).py}. Explain
  lines 25--26: how do \(W1 = Z1\) and
  \texttt{W2\ =\ rho*Z1\ +\ np.sqrt(1-rho**2)*Z2} produce two Brownians
  with \(dW1 dW2 = ρ dt\)? Why is this the Cholesky factor of a 2×2
  correlation matrix? \textbf{A.} \(Z1, Z2\) are i.i.d. N(0,1). Then
  \(Var(W2) = ρ²·1 + (1−ρ²)·1 = 1\), and
  \(Cov(W1,W2) = E[Z1(ρZ1 + √(1−ρ²)Z2)] = ρ\), so \((W1,W2)\) are
  standard normals with correlation ρ. The map \([Z1,Z2] ↦ [W1,W2]\)
  uses the lower-triangular matrix \([[1,0],[ρ,√(1−ρ²)]]\), which is
  exactly the Cholesky factor \(L\) of the correlation matrix
  \([[1,ρ],[ρ,1]] = L Lᵀ\). Scaling by √T (or working in increments)
  gives Brownians with \(dW1 dW2 = ρ dt\).
\item
  Compare line 29 (\(drift r − 0.5σ²\), under Q) with line 65
  (\(drift r + 0.5σ1²\), under QS1). Why does the sign in front of the
  \(½σ1²\) term flip, and which result from Exercise 3(c) does that come
  from? \textbf{A.} Under Q the asset has log-drift \(r − ½σ1²\) (Itô
  correction for the lognormal). Under the measure QS1 with numéraire
  S1, by Girsanov the Brownian driving S1 picks up a \(+σ1 dt\) drift
  shift, which turns S1's log-drift into \(r − ½σ1² + σ1² = r + ½σ1²\).
  That sign flip is exactly the change-of-numéraire / Radon--Nikodym
  result derived in 3(c): changing numéraire from the bank account to S1
  adds \(σ1²T\) to S1's exponent (and \(ρσ1σ2 T\) to S2's, line 66).
\item
  Line 71 returns \texttt{V\ =\ S0{[}0{]}*np.maximum(AT-K,0)/S1} with
  \textbf{no} \(e^{−rT}\) factor, whereas line 35 (under Q)
  \textbf{does} discount. Derive why the discount factor is absent under
  QS1. \textbf{A.} Under Q,
  \texttt{V0\ =\ e\^{}\{−rT\}\ E\^{}Q{[}payoff{]}} --- the bank account
  \(e^{rt}\) is the numéraire, hence the discount. Under QS1 the
  numéraire is S1 itself, so the pricing formula is
  \texttt{V0\ =\ S1(0)\ E\^{}\{QS1\}{[}payoff\ /\ S1(T){]}}. The
  \texttt{1/S1} on line 71 is the division by the terminal numéraire and
  the \(S0[0]\) prefactor is \(S1(0)\); there is no \(e^{−rT}\) because
  the deterministic bank-account discount has been replaced by the
  stochastic numéraire S1. The two estimators target the same V0.
\item
  What exactly do lines 38--39 (and 74--75) return, and why is
  \(ddof=1\) used in \texttt{np.std}? What would change if you used
  \(ddof=0\)? \textbf{A.} They return \(V_aprox = mean(V)\) (the MC
  price estimator) and \texttt{SE\_aprox\ =\ std(V,\ ddof=1)/√M} (its
  standard error). \(ddof=1\) gives the unbiased sample variance (divide
  by \(M−1\)), which is the correct estimator of the population variance
  feeding the CLT-based SE. \(ddof=0\) divides by \(M\), slightly
  underestimating the variance and hence the SE; for the large M here
  (10k--100k) the difference is negligible (\textasciitilde factor
  √(M/(M−1)) ≈ 1), but ddof=1 is the principled choice.
\item
  Lines 118--122 construct the \texttt{1/√M} reference curves. What is
  the constant \texttt{c\_Q\ =\ SEQ{[}0{]}*np.sqrt(values\_M{[}0{]})}
  doing, and why does plotting \texttt{c\_Q/np.sqrt(M)} test the CLT
  rate rather than assume it? \textbf{A.} \(c_Q = SE(M0)·√M0\) recovers
  the proportionality constant (≈ the population std of the payoff) from
  the \emph{first} data point only, anchoring the reference line
  \texttt{c\_Q/√M}. Because the constant is fixed from one M and the
  rest of the curve is pure \texttt{1/√M}, overlaying it on the measured
  SEs is a genuine test: if the measured SEs fall on the line, the
  empirical decay rate really is −1/2. It tests the rate because only
  the level (not the slope) was fitted.
\item
  \texttt{np.random.seed(44)} is set on line 14 (Q) and line 50 (QS1).
  What is the consequence of reseeding identically in both functions for
  the \emph{fairness} of the Q-vs-QS1 SE comparison? Is the convergence
  \texttt{O(1/√M)} either way? \textbf{A.} Reseeding with the same 44
  means both measures draw the \emph{same} underlying \(Z1, Z2\)
  streams, so the SE comparison is on common random numbers --- the
  difference between the two SEs is attributable to the change of
  measure (the variance of the payoff functional), not to lucky/unlucky
  draws. That makes the comparison fair (paired). The \texttt{O(1/√M)}
  rate holds under either measure regardless --- it is the CLT and
  depends only on finite payoff variance --- what changes is the
  \emph{constant} multiplying \texttt{1/√M}.
\item
  Using Scenario-1 numbers (\(σ1=σ2=0.2\), \(ρ=0.9999\)), compute
  \( = √(σ1²+σ2²−2ρσ1σ2)\) and explain why this near-zero value is the
  reason the QS1 SE on line 75 is so small. \textbf{A.}
  \(² = 0.04 + 0.04 − 2(0.9999)(0.2)(0.2) = 0.08 − 0.079992 = 8e-6\),
  so \( ≈ 0.00283\). With ρ≈1 the two assets move almost identically,
  so the basket spread \(S1−S2\) (and the payoff's randomness after
  switching numéraire to S1) is governed by this tiny \(\). Under QS1
  the estimator's variance scales with \(²\), so it nearly collapses to
  a deterministic value --- the payoff has almost no variance --- making
  the QS1 SE on line 75 orders of magnitude smaller than the Q SE. This
  is the variance-reduction win of the S1-numéraire formulation for
  highly correlated baskets.
\end{enumerate}

\hypertarget{assignment-set-1}{%
\section{Assignment set 1}\label{assignment-set-1}}

Open the named file and answer at the line/function level. Point to the
exact line, don't paraphrase the report.

\begin{enumerate}
\def\labelenumi{\arabic{enumi}.}
\item
  \textbf{\texttt{2bCF\ (2).py:26}} ---
  \texttt{fk\ =\ (2/(b-a))\ *\ Im\{\ char(wk)\ *\ exp(-i\ wk\ a)\ \}}.
  Explain term by term where this comes from. Why the imaginary part?
  Why the \texttt{2/(b-a)} prefactor (and not \texttt{1/(b-a)})? Why
  does \(k\) start at 1 (\texttt{:22}) with no \(k=0\) term? \textbf{A.}
  This is the Fourier-\emph{sine} coefficient. Expanding \(f\) in
  \texttt{\{sin(kπ(x−a)/(b−a))\}} on \([a,b]\), the coefficient is
  \texttt{f\_k\ =\ (2/(b−a))\ ∫\_a\^{}b\ f(x)\ sin(ω\_k(x−a))\ dx} with
  \texttt{ω\_k\ =\ kπ/(b−a)}. Writing
  \texttt{sin(ω\_k(x−a))\ =\ Im\{e\^{}\{iω\_k(x−a)\}\}} and recognising
  \texttt{∫\ f(x)\ e\^{}\{iω\_k\ x\}\ dx\ =\ φ(ω\_k)} (the
  characteristic function), the integral becomes
  \texttt{Im\{φ(ω\_k)\ e\^{}\{−iω\_k\ a\}\}}. So the \(Im\) is because
  the basis is sine (the odd part), \(char(wk)=φ(ω_k)\) is the CF, and
  \(exp(−i ω_k a)\) is the phase shift to the \([a,b]\) origin. The
  \texttt{2/(b−a)} is the standard L²-normalisation of sine series (not
  \texttt{1/(b−a)}, which is the cosine \(k=0\) average term). \(k\)
  starts at 1 because \(sin(0)=0\), so there is no \(k=0\) sine mode.
\item
  \textbf{\texttt{2bCF\ (2).py:30-31}} --- the reconstruction loop.
  Write out the sum being evaluated and say why it approximates \(f(x)\)
  on \([a,b]\). What sets the accuracy: \(N\), or the range
  \([a,b]=[-10,10]\) (\texttt{:4-5})? \textbf{A.} It evaluates
  \texttt{f\_recovered(x)\ =\ Σ\_\{k=1\}\^{}N\ f\_k\ sin(kπ(x−a)/(b−a))},
  the truncated Fourier-sine partial sum. It approximates \(f\) because
  \texttt{\{sin(kπ(x−a)/(b−a))\}} is an orthogonal basis of \(L²[a,b]\)
  and \(f_k\) are the projections; truncating at N keeps the N lowest
  modes. Both control accuracy but differently: the \emph{truncation}
  \([a,b]\) must be wide enough to contain essentially all the mass of
  \(f\) (so the series error from the tails is negligible --- here
  \([−10,10]\) is \textasciitilde10σ for N(0,1), more than enough), and
  given a good range, increasing \(N\) drives the series-truncation
  error down (exponentially for this smooth density). With \([−10,10]\)
  fixed and generous, \(N\) is the binding accuracy knob.
\item
  \textbf{\texttt{2dCF\ (2).py} --- function \(phi_BS\)} --- it returns
  \texttt{exp(i\ w\ mu\_y\ -\ 0.5\ var\_y\ w\^{}2)} with
  \(mu_y = x + (r - 0.5 sigma^2) T\). Why is this the characteristic
  function of \texttt{y\ =\ ln(S\_T/K)} and not of \(ln(S_T)\)? Where
  does \texttt{x\ =\ x0\ =\ log(S0/K)} enter? \textbf{A.} Under Q,
  \(ln(S_T) = ln(S0) + (r − ½σ²)T + σ√T Z\) is normal with mean
  \(ln(S0)+(r−½σ²)T\) and variance \(σ²T\). Define
  \texttt{y\ =\ ln(S\_T/K)\ =\ ln(S\_T)\ −\ ln\ K}; this just shifts the
  mean by \(−ln K\), giving mean
  \texttt{ln(S0/K)\ +\ (r−½σ²)T\ =\ x0\ +\ (r−½σ²)T}. The CF of a normal
  \(N(μ_y, σ²T)\) is \(exp(iωμ_y − ½σ²T ω²)\), exactly the return value.
  So it is the CF of \texttt{y\ =\ ln(S\_T/K)} (centred on the strike)
  precisely because \(mu_y\) carries \texttt{x\ =\ x0\ =\ ln(S0/K)};
  that \(x\) is the only place the spot/strike ratio enters, and it is
  what makes the expansion be of the log-\emph{moneyness} rather than
  the raw log-price.
\item
  \textbf{\texttt{2dCF\ (2).py} --- cumulant range block} ---
  \(a = c1 - L*sqrt(c2 + sqrt(c4))\) with \(c1 = x0 + mu*T\),
  \(c2 = T sigma^2\), \(c4 = 0\), \(L = 10\). Derive \(c1\) and \(c2\)
  as the first and second cumulants of the log-price. Why is \(c4 = 0\)
  legitimate for Black--Scholes? What breaks if \(L\) is too small?
  \textbf{A.} For
  \texttt{y\ =\ ln(S\_T/K)\ \textasciitilde{}\ N(x0+(r−½σ²)T,\ σ²T)},
  the first cumulant (mean) is \(c1 = x0 + μT\) with \(μ = r−½σ²\)
  (line: \(mu = (r−0.5σ²)\)), and the second cumulant (variance) is
  \(c2 = σ²T\). A normal distribution has \emph{all} cumulants of order
  ≥3 equal to zero, so \(c4 = 0\) is exact for Black--Scholes --- that
  is legitimate, not an approximation. The Fang--Oosterlee rule
  \([c1 ± L√(c2+√c4)]\) then reduces to \(c1 ± Lσ√T\), i.e.~±10 standard
  deviations. If \(L\) is too small the interval \([a,b]\) does not
  contain the effective support of the density/payoff, the truncated
  sine/cosine series misses mass, and the COS price acquires a large,
  non-decaying \emph{truncation} error that no increase in \(N\) can
  fix.
\item
  \textbf{\texttt{2dCF\ (2).py} --- functions \(phi_k\) and \(chi_k\)}
  --- note the naming clash: the function called \(phi_k\) is the
  report's \(psi_k\) (integral of \(sin\)), and \(chi_k\) matches the
  report's \(chi_k\) (integral of \(e^y sin\)). Reconcile the code with
  your 2c derivation and confirm \texttt{V\_k\_put} pairs them
  correctly. Why does the function name \(phi_k\) collide awkwardly with
  the characteristic function \(phi\)? \textbf{A.} The code's \(phi_k\)
  computes \texttt{(1/ω\_k){[}cos(ω\_k(c−a))\ −\ cos(ω\_k(d−a)){]}},
  which is \texttt{∫\_c\^{}d\ sin(ω\_k(y−a))\ dy} --- that is the
  report's \(ψ_k\) (the bare-sine integral). The code's \(chi_k\)
  computes \texttt{∫\_c\^{}d\ e\^{}y\ sin(ω\_k(y−a))\ dy} via the
  standard \texttt{e\^{}y(sin\ −\ ω\_k\ cos)/(1+ω\_k²)} antiderivative
  --- that matches the report's \(χ_k\). For a put, intrinsic
  \(K − S = K(1 − e^y)\) on \(y<0\), so
  \texttt{V\_k\ ∝\ ∫(K\ −\ Ke\^{}y)sin(...)\ =\ K(ψ\_k\ −\ χ\_k)};
  \texttt{V\_k\_put} returns \texttt{(2K/(b−a))(−χ\_k\ +\ ψ\_k)},
  i.e.~\texttt{K(ψ\_k\ −\ χ\_k)} --- correctly paired. The name
  \(phi_k\) is unfortunate because \(phi\)/\(phi_BS\) already denote the
  characteristic function; the local \(phi_k\) is a payoff-integral,
  conceptually unrelated, so the clash is purely cosmetic naming, not a
  bug. \emph{(confirm this matches your own reasoning/observation)}
\item
  \textbf{\texttt{2dCF\ (2).py} --- \texttt{V\_k\_put}} --- returns
  \texttt{(2K/(b-a))\ *\ (-chi\_k(...,a,0)\ +\ phi\_k(...,a,0))}. Why
  are the integration limits \((a, 0)\) and not \((a, b)\)? Tie this to
  the put payoff being zero for \(y > 0\). \textbf{A.} The put payoff in
  log-moneyness is \(K max(1 − e^y, 0)\), which is positive only for
  \(y < 0\) (i.e.~\(S_T < K\)) and identically zero for \(y > 0\). So
  the coefficient integral
  \texttt{∫\_a\^{}b\ payoff·sin\ =\ ∫\_a\^{}0\ K(1−e\^{}y)sin} --- the
  \([0,b]\) part contributes nothing. The limits \((a, 0)\) simply
  restrict to the region where the put has value; integrating over
  \((a,b)\) would add a zero contribution but the analytic
  \texttt{ψ\_k,\ χ\_k} are written for the active region, so \(0\) is
  the correct upper limit.
\item
  \textbf{\texttt{2dCF\ (2).py} --- \texttt{sin\_put\_price}} --- final
  line is \(exp(-rT) * dot(cf_part, Vk)\). Why \(e^{-rT}\) and not
  \(e^{-r Delta t}\) with a different \(Delta t\)? Why is the price a
  dot product of the CF part with the payoff coefficients (what theorem
  makes the product of two functions become a product of their sine
  coefficients)? \textbf{A.} It is a one-shot European option to
  maturity T with no intermediate steps, so the discount is over the
  whole horizon, \(e^{−rT}\); there is no time grid hence no \(Δt\). The
  price is
  \texttt{e\^{}\{−rT\}E{[}payoff{]}\ =\ e\^{}\{−rT\}∫\ payoff(y)\ f(y)\ dy}.
  Expanding both \(f\) and the payoff in the same orthogonal sine basis
  on \([a,b]\) and applying Parseval/Plancherel, the integral of the
  product equals the (weighted) sum of products of their coefficients:
  \texttt{∫\ f·payoff\ ≈\ Σ\_k\ Im\{φ(ω\_k)e\^{}\{−iω\_k\ a\}\}\ V\_k}.
  That is exactly \(dot(cf_part, Vk)\) --- Parseval's theorem turns the
  function inner product into a coefficient dot product.
\item
  \textbf{\texttt{2eCF\ (2).py:46}} ---
  \texttt{return\ (2/pi)\ *\ (im/ks)\ @\ (1\ -\ cosine).T}. Derive this
  from integrating the sine density term by term. Where do the
  \texttt{1/k} factor and the \texttt{2/pi} prefactor come from, and why
  is the integrated \(sin\) term \((1 - cos(...))\)? \textbf{A.} The CDF
  is \(F(z) = ∫_a^z f(x)dx\). With
  \texttt{f(x)\ ≈\ Σ\_k\ f\_k\ sin(ω\_k(x−a))} and
  \texttt{f\_k\ =\ (2/(b−a))\ Im\{φ(ω\_k)e\^{}\{−iω\_k\ a\}\}} (= \(im\)
  in code), integrate each mode:
  \texttt{∫\_a\^{}z\ sin(ω\_k(x−a))dx\ =\ (1/ω\_k){[}1\ −\ cos(ω\_k(z−a)){]}}.
  The lower-limit \(cos(0)=1\) produces the \(1\) in \((1 − cos(...))\).
  Now \texttt{(2/(b−a))·(1/ω\_k)\ =\ (2/(b−a))·(b−a)/(kπ)\ =\ 2/(kπ)},
  so the \((b−a)\) cancels and leaves the \texttt{2/π} prefactor and the
  \texttt{1/k} factor (\texttt{im/ks}). Hence
  \texttt{F(z)\ =\ (2/π)\ Σ\_k\ (im\_k/k)(1\ −\ cos(kπ(z−a)/(b−a)))},
  exactly the matrix expression \texttt{(2/π)(im/ks)\ @\ (1−cosine).T}.
\item
  \textbf{\texttt{2eCF\ (2).py:14-15}} --- fixed range \(a=-2, b=2\).
  For \texttt{ln(S\_T/S\_0)\ \textasciitilde{}\ N(mu,\ sigma\_aux\^{}2)}
  with \(mu ≈ -0.015\), \(sigma_aux = 0.3\), how many standard
  deviations is \([-2,2]\)? Is that wide enough? Why does the 2e error
  plateau at \(3.67e-11\) rather than reaching \(1e-16\) as 2d does?
  \textbf{A.} Here \(σ_aux = σ√T = 0.3·√1 = 0.3\) and
  \(μ ≈ (r−½σ²)T = 0.03−0.045 = −0.015\). The interval reaches
  \texttt{(2\ −\ (−0.015))/0.3\ ≈\ 6.7} σ above and
  \texttt{(−2\ −\ (−0.015))/0.3\ ≈\ −6.6} σ below the mean --- roughly
  ±6.7σ, which is wide enough that the Gaussian tail mass outside is
  \textasciitilde10⁻¹¹. That tail mass is precisely the floor: the CDF
  reconstruction can be no more accurate than the density mass it
  ignores beyond \([a,b]\), so 2e plateaus at ≈\(3.67e-11\) (the
  truncation tail), whereas 2d's \emph{price} is an integral against a
  payoff that decays, letting it reach machine precision. The plateau is
  a domain-truncation error, not a series-truncation error ---
  increasing \(N\) will not push it below the tail. \emph{(confirm this
  matches your own observation)}
\item
  \textbf{\texttt{3aCF\ (2).py:26}} ---
  \texttt{S\_T\ =\ S\_0\ *\ exp(drift*T\ +\ sigma1*Wt{[}:,None{]}\ +\ sigma2*Bi)}
  with \texttt{drift\ =\ r\ -\ 0.5*(sigma1\^{}2\ +\ sigma2\^{}2)}.
  Derive this exact terminal law from the SDE (your 3a Ito derivation).
  Why is \texttt{0.5*(sigma1\^{}2\ +\ sigma2\^{}2)} the right convexity
  correction, and why is no time-stepping needed? \textbf{A.} The asset
  follows \texttt{dS/S\ =\ r\ dt\ +\ σ1\ dW\ +\ σ2\ dB} with \(W\) the
  common factor and \(B\) the idiosyncratic factor, independent.
  Applying Itô to \(ln S\):
  \(d ln S = (r − ½(σ1²+σ2²))dt + σ1 dW + σ2 dB\), since the total
  quadratic variation of the log is \((σ1²+σ2²)dt\) (W and B
  independent, so no cross term). Integrating to T gives
  \texttt{ln(S\_T/S0)\ =\ (r\ −\ ½(σ1²+σ2²))T\ +\ σ1\ W\_T\ +\ σ2\ B\_T}
  with \texttt{W\_T,\ B\_T\ \textasciitilde{}\ N(0,T)}. So
  \(½(σ1²+σ2²)\) is exactly the Itô convexity correction for the
  \emph{combined} variance. Because the drift and diffusion are
  constant, the SDE has a closed-form lognormal solution at T --- Itô
  integrates exactly --- so no Euler time-stepping is needed; one draw
  of \texttt{(W\_T,\ B\_i)} gives \(S_T\) exactly. (Note:
  \texttt{Wt{[}:,None{]}} broadcasts the single common shock across all
  1000 stocks; \(Bi\) is the per-stock idiosyncratic shock.)
\item
  \textbf{\texttt{3aCF\ (2).py:42}} ---
  \texttt{q\ =\ np.quantile(L\_T,\ 1\ -\ alpha)} with
  \(alpha = 1 - 0.95\). Given the assignment's definition
  \texttt{P(L\_T\ \textgreater{}=\ q\_alpha)\ \textless{}=\ alpha},
  justify taking the \(0.95\) empirical quantile. What interpolation
  does \texttt{np.quantile} use, and could that matter for a discrete
  \(L_T\)? \textbf{A.} \(alpha = 1 − 0.95 = 0.05\), so
  \(1 − alpha = 0.95\) and \texttt{np.quantile(L\_T,\ 0.95)} returns the
  value \(q\) with \(P(L_T ≤ q) ≈ 0.95\), equivalently
  \(P(L_T > q) ≈ 0.05 = alpha\) --- the 95\% VaR-style quantile
  required. \texttt{np.quantile} defaults to \emph{linear} interpolation
  between order statistics. For a discrete integer-valued \(L_T\) (it is
  a count, \(Σ 1{S_i>K}\)) the true quantile is an integer, and linear
  interpolation can return a non-integer between two adjacent counts;
  this introduces a small interpolation artefact, though with 500k paths
  it is negligible. A \(method='lower'\)/\('higher'\) choice would
  respect the \(≤ alpha\) inequality strictly. \emph{(confirm this
  matches your own reasoning)}
\item
  \textbf{\texttt{3bCF\ (2).py:53}} ---
  \texttt{ratio\ =\ exp(-(mu\_shift/T)*W\_q\ +\ mu\_shiftsq)} with
  \texttt{mu\_shiftsq\ =\ mu\_shift\^{}2/(2T)}. Derive this \texttt{p/q}
  from two Gaussians \(p~N(0,T)\), \(q~N(1.5,T)\) and confirm it equals
  \(exp(-1.5 w + 1.125)\) at \(T=1\). Why is only \(W\) reweighted and
  not the \(B_i\) (\texttt{:44})? \textbf{A.} The likelihood ratio is
  \texttt{p(w)/q(w)\ =\ exp(−w²/2T)\ /\ exp(−(w−μ)²/2T)\ =\ exp({[}−w²\ +\ (w−μ)²{]}/2T)\ =\ exp((−2μw\ +\ μ²)/2T)\ =\ exp(−(μ/T)w\ +\ μ²/(2T))}
  with \(μ = mu_shift = 1.5\). That is exactly
  \texttt{exp(−(mu\_shift/T)W\_q\ +\ mu\_shiftsq)}. At \(T=1\),
  \texttt{μ²/(2T)\ =\ 2.25/2\ =\ 1.125}, so
  \(ratio = exp(−1.5 w + 1.125)\), confirming the comment. Only \(W\) is
  reweighted because the importance shift was applied only to the
  \emph{common} factor \(W\) (sampled from \(N(1.5,T)\) on line 43); the
  idiosyncratic \(B_i\) are still drawn from their original \(N(0,T)\)
  (line 44), so their density is unchanged and contributes a factor of 1
  to \texttt{p/q}. Shifting only \(W\) is enough because the rare event
  (many \(S_i > K\)) is driven by the common factor.
\item
  \textbf{\texttt{3cCF\ (2).py:54}} ---
  \texttt{ratio\ =\ sqrt(2)*exp((-W\_q\^{}2\ -\ 3\ W\_q\ +\ 2.25)/(4T))}.
  Derive this from \(p~N(0,T)\), \(q~N(1.5,2T)\). Where does the
  \texttt{sqrt(2)} prefactor come from? Explain why this mean+var shift
  is \emph{less} efficient for the 95\% quantile than the pure mean
  shift in 3b. \textbf{A.}
  \texttt{p(w)\ =\ (2πT)\^{}\{−1/2\}\ exp(−w²/2T)},
  \texttt{q(w)\ =\ (2π·2T)\^{}\{−1/2\}\ exp(−(w−1.5)²/(2·2T))}. Then
  \texttt{p/q\ =\ √(2T/T)·exp(−w²/2T\ +\ (w−1.5)²/4T)\ =\ √2·exp({[}−2w²\ +\ (w−1.5)²{]}/4T)\ =\ √2·exp((−w²\ −\ 3w\ +\ 2.25)/4T)},
  exactly the code. The \(√2\) is the ratio of the normalising constants
  \texttt{√(var\_q/var\_p)\ =\ √(2T/T)}. It is less efficient because
  inflating the variance (\(2T\) vs \(T\)) makes the likelihood ratio
  \texttt{p/q} much more dispersed --- its tail produces a few paths
  with very large weights, raising the \emph{variance of the weights}
  and hence of the weighted quantile estimator. The pure mean shift of
  3b keeps the variance matched, so the weights are bounded and
  well-behaved; for a quantile (not a mean) of a common-factor-driven
  rare event, the mean shift alone already moves mass into the tail
  without the weight blow-up, so it wins. \emph{(confirm this matches
  your own observation)}
\item
  \textbf{\texttt{3bCF\ (2).py:20-34} / \texttt{3cCF\ (2).py:20-34}} ---
  \(weighted_quantile\): \texttt{cumsum(w\_sorted)/sum(w\_sorted)} then
  \(searchsorted(cdf, q)\). Why must you sort by \(values\) first and
  carry the weights along? Why does 3a use \texttt{np.quantile} but
  3b/3c need this custom weighted version --- and does that estimator
  mismatch bias the efficiency comparison against 3a? \textbf{A.} A
  quantile is read off the CDF, and the empirical CDF is
  \(cumsum of weights ordered by the value axis\); so you must
  \(argsort(values)\) and permute the weights by the \emph{same}
  \(sorter\) (lines 25--27) to keep each weight attached to its
  observation. \texttt{cumsum(w\_sorted)/sum(w\_sorted)} is then the
  weighted empirical CDF and \(searchsorted(cdf, q)\) finds the first
  value whose cumulative mass reaches \(q\). 3a is plain MC under the
  true measure \(p\), where every path has equal weight \texttt{1/N}, so
  \texttt{np.quantile} (equal-weight) is correct. 3b/3c sample under
  \(q\), so each path carries its likelihood ratio \texttt{p/q} as a
  weight, and the unweighted \texttt{np.quantile} would be biased ---
  the weighted version restores an unbiased estimate of the
  \(p\)-quantile. The estimator mismatch does \emph{not} bias the
  comparison: both target the same population quantile of \(L_T\) under
  \(p\); what differs is the variance of the estimator, which is exactly
  the efficiency being compared. (One caveat: with linear interpolation
  in 3a vs \(searchsorted\) index lookup in 3b/3c, the tiny
  discretisation conventions differ, but at 500k paths this is
  immaterial.) \emph{(confirm this matches your own reasoning)}
\end{enumerate}

\hypertarget{exercise-set-2}{%
\section{Exercise set 2}\label{exercise-set-2}}

Files: - \texttt{assignments/exercise-set-2/exercise1\_def\ (2).py} (LSM
/ Longstaff--Schwartz) -
\texttt{assignments/exercise-set-2/exercise2\_def\ (2).py} (FTCS finite
differences)

\hypertarget{exercise-1-lsm}{%
\subsection{Exercise 1 --- LSM}\label{exercise-1-lsm}}

\begin{enumerate}
\def\labelenumi{\arabic{enumi}.}
\item
  \textbf{\texttt{exercise1\_def\ (2).py:96-127}} --- walk the backward
  LSM loop out loud. Which line builds the discounted continuation
  target \(Y\), which applies the ITM mask, which fits the regression,
  which makes the exercise decision, and which zeroes the later cash
  flows? \textbf{A.} The loop runs \(i\) from \(T−2\) down to \(0\).
  Line 101 \(Y = C[:, i+1:] @ discount_factors\) builds the discounted
  continuation target (all realised future cash flows discounted back to
  \(i+1\)). Line 102 \(itm_indexes = S[:,i+1] < K\) is the ITM mask,
  applied on lines 103--104 to get \texttt{X\_itm,\ Y\_itm}. Lines
  107--108 fit the degree-2 regression
  (\texttt{model.fit(X\_itm,\ Y\_itm)}). The exercise decision is lines
  122--124:
  \texttt{if\ exercise{[}j{]}\ \textgreater{}\ continuation{[}j{]}} set
  \(C[j,i] = exercise[j]\). Line 125 \(C[j,i+1:] = 0.00\) zeroes the
  later cash flows on exercised paths.
\item
  \textbf{\texttt{:99-101}} --- \(Y\) is the sum of all future cash
  flows discounted by \(exp(-r*dt*(u-i))\). Why discount to time \(i+1\)
  (not to \(0\))? What would change in the comparison at
  \texttt{:117-120} if you discounted to \(0\) instead? \textbf{A.}
  \(u = arange(i+1, ncols)\) and the factor is \(exp(−r·dt·(u−i))\),
  which discounts each future cash flow to time \(i\) (one step ahead of
  \(i\)). The exercise decision at step \(i+1\) compares the
  \emph{intrinsic value at \(i+1\)} against the \emph{continuation value
  at \(i+1\)}, so both sides must be expressed at the same date. The
  regression target \(Y\) and the exercise payoff
  \(exercise = K − S[:,i+1]\) must live at the same time point;
  discounting \(Y\) to 0 instead would put continuation at time 0 while
  intrinsic stays at \(i+1\), making the comparison on lines 117--124
  mix two time points and corrupt the early-exercise decision. (The
  absolute level only needs to be consistent on both sides of the
  inequality.) \emph{(confirm this matches your own reasoning)}
\item
  \textbf{\texttt{:102-104}} --- the ITM mask \(S[:,i+1] < K\). Why does
  Longstaff--Schwartz regress \emph{only} on in-the-money paths? What
  bias appears if you include OTM paths? \textbf{A.} Early exercise is
  only ever relevant where the option is in the money (a put with
  \(S ≥ K\) has zero intrinsic, so you would never exercise).
  Restricting the regression to ITM paths concentrates the basis
  functions on the region where the exercise boundary actually lives,
  giving a far better local fit of the continuation value near the
  boundary. Including OTM paths forces the low-degree polynomial to also
  fit the flat zero-intrinsic region, distorting the fitted continuation
  function precisely near the boundary, which biases the exercise
  decision and typically \emph{underprices} the option. This is the
  original Longstaff--Schwartz (2001) prescription.
\item
  \textbf{\texttt{:107-112}} --- degree-2 monomial pipeline; you
  reported \(E[Y|X] = -1.06999 + 2.98341 X - 1.81358 X^2\) at \(t=2\).
  Reproduce how \(a0,a1,a2\) are fit and what \(X\), \(Y\) are at that
  step. \textbf{A.} At \(i=1\) (so step \(t=i+1=2\)), \(X\) is
  \(S[:,2]\) masked to ITM paths (line 98 builds \(X = S[:,i+1]·1{S<K}\)
  reshaped to a column; \(X_itm\) keeps the ITM rows), and \(Y\) is
  \(Y_itm\), the discounted realised cash flows from \(t=3\) back to
  \(t=2\) (\(C[:,2:] @ exp(−r·dt)\)). The pipeline
  \texttt{PolynomialFeatures(degree=2,\ include\_bias=False)} builds
  columns \([X, X²]\), and
  \texttt{LinearRegression(fit\_intercept=True)} does ordinary least
  squares, returning \(intercept_ = a0\) and \(coef_ = (a1, a2)\) (lines
  111--112). Those OLS coefficients are the reported
  \(−1.06999 + 2.98341 X − 1.81358 X²\), the conditional-expectation
  estimate \(Ê[Y|X]\). \emph{(confirm the numbers match your printout)}
\item
  \textbf{\texttt{:122-127}} --- the exercise rule sets \(C[j,i]\) to
  the intrinsic value and zeroes \(C[j,i+1:]\). Why must all later cash
  flows be zeroed once you exercise? What does the stopping matrix at
  \texttt{:131} then represent? \textbf{A.} An American option is
  exercised at most \emph{once}; once you exercise at \(i\), the
  contract is dead and there can be no later cash flow on that path.
  Zeroing \(C[j,i+1:]\) enforces a single realised cash flow per path,
  so summing the row gives the path's payoff without double counting.
  The stopping matrix at line 131 (\texttt{np.where(C!=0,1,0)}) then has
  exactly one \(1\) per exercised path marking its optimal stopping time
  --- it is the discrete exercise/stopping policy.
\item
  \textbf{\texttt{:147-156}} --- American vs European price. Why does
  the European price use only the terminal column \(C_european\) while
  the American price sums each path's single realized cash flow?
  \textbf{A.} The European option can only be exercised at \(T\), so its
  cash flow is the terminal payoff
  \texttt{C\_european\ =\ max(K−S\_T,0)} (saved on line 93 before the
  loop overwrites \(C\)), discounted once over the full horizon:
  \(mean(C_european·exp(−r·dt·T))\) (line 156). The American option may
  exercise at any step, and after the backward loop each path's row of
  \(C\) holds its single realised cash flow at its optimal time; line
  150 discounts each by \(exp(−r·dt·t)\) for its own \(t\) and sums per
  path, then averages (lines 150--151). Different treatment because the
  two have different exercise rights.
\item
  \textbf{\texttt{:177-178}} --- antithetic paths (\(Z\) then \(-Z\)).
  State the covariance argument for why this reduces variance for this
  payoff. Does the s.e. \texttt{std(cashflow)/sqrt(N)} at \texttt{:258}
  remain correct when the \(N\) draws are antithetically paired (not
  independent)? \textbf{A.} With antithetics the estimator averages
  \(½(f(Z)+f(−Z))\). Its variance is
  \(¼[Var f(Z) + Var f(−Z) + 2Cov(f(Z),f(−Z))] = ½ Var f(Z)(1 + ρ)\)
  with \(ρ = Corr(f(Z),f(−Z))\). For a monotone payoff like a put
  (decreasing in \(S\), hence in \(Z\)), \(f(Z)\) and \(f(−Z)\) are
  negatively correlated, \(ρ<0\), so the bracket is below \(Var f(Z)\):
  variance is reduced. The caveat on line 258: \texttt{std(cashflow)/√N}
  treats the \(N\) draws as i.i.d., but the antithetic pairs are
  \emph{dependent}, so this formula is biased (it ignores the
  within-pair negative covariance and typically \emph{overstates} the
  true SE). The correct SE computes the std over the \texttt{N/2}
  \emph{pair averages} divided by \texttt{√(N/2)}. So the reported s.e.
  is conservative/incorrect in form. \emph{(confirm --- this is a real
  subtlety Fang may push on)}
\item
  \textbf{\texttt{:185-191}} --- \(laguerre_basis\) uses
  \texttt{x\ =\ S/K} and weights \texttt{exp(-x/2)}. Why normalise by
  \(K\), and what numerical problem (underflow) does \texttt{exp(-x/2)}
  with \texttt{x=S/K} avoid that raw monomials in \(S\) would hit for
  \(S∈[36,44]\)? \textbf{A.} Normalising \texttt{x\ =\ S/K} makes the
  regressor O(1) (around 0.9--1.1 for \(S≈36–44\), \(K=40\)) instead of
  O(40); raw monomials \(S, S², S³\) in absolute price would span
  \textasciitilde36 to \textasciitilde85k, producing a badly conditioned
  (near-collinear, huge-dynamic-range) design matrix and ill-conditioned
  normal equations. The Laguerre weight \texttt{exp(−x/2)} with the
  small normalised \(x\) keeps the basis bounded and well-scaled; with
  raw \(S\) the weight \texttt{exp(−S/2)} would underflow to
  \textasciitilde{}\(exp(−18)≈10⁻⁸\) (or smaller), collapsing the basis
  to numerical zero. Normalising by \(K\) is exactly the rescaling
  Longstaff--Schwartz use to keep the regression well-conditioned.
\item
  \textbf{\texttt{:257} vs \texttt{exercise2\_def\ (2).py:223}} --- one
  uses \(abs(american-european)\), the other the \emph{signed}
  difference for the early-exercise value. Which matches the paper's
  ``early-exercise value'' column, and can LSM-American ever dip below
  European by Monte-Carlo noise? \textbf{A.} The signed difference
  \(american − european\) (\texttt{exercise2\_def\ (2).py:223}) matches
  the paper's ``early-exercise premium'' column --- early-exercise value
  is by definition \(American − European ≥ 0\) and its sign carries
  information. The \(abs(...)\) on \texttt{exercise1\_def\ (2).py:257}
  discards that sign. The signed form can go slightly negative: LSM is a
  \emph{suboptimal} (low-biased) estimator of the American price (the
  regression-based exercise rule is never better than optimal, and any
  policy is suboptimal in-sample noise aside), and on top of that
  Monte-Carlo noise can push the simulated American below the
  closed-form European, giving a small negative premium. The \(abs\)
  hides exactly this diagnostic, so the signed difference is the honest
  column. \emph{(confirm this matches what you observed in the
  deep-OTM/short-T rows)}
\end{enumerate}

\hypertarget{exercise-2-ftcs}{%
\subsection{Exercise 2 --- FTCS}\label{exercise-2-ftcs}}

\begin{enumerate}
\def\labelenumi{\arabic{enumi}.}
\setcounter{enumi}{9}
\item
  \textbf{\texttt{exercise2\_def\ (2).py:35-40}} --- derive the explicit
  evolution matrix
  \texttt{F\ =\ (1\ -\ r\ k)\ I\ +\ 0.5\ k\ σ²\ /\ h²\ ·\ T2\ +\ k\ (r\ -\ 0.5\ σ²)/(2h)\ ·\ T1}
  from the Black--Scholes PDE written in \(m = log S\). Identify which
  term is the reaction, which the diffusion, which the advection.
  \textbf{A.} Writing \(m = ln S\), the BS PDE in
  \emph{time-to-maturity} form is
  \texttt{∂U/∂τ\ =\ ½σ²\ ∂²U/∂m²\ +\ (r\ −\ ½σ²)\ ∂U/∂m\ −\ rU}.
  Discretise with forward Euler in time (step \(k\)) and central
  differences in \(m\) (step \(h\)):
  \texttt{(U\^{}\{n+1\}\ −\ U\^{}n)/k\ =\ ½σ²/h²\ (T2\ U\^{}n)\ +\ (r−½σ²)/(2h)(T1\ U\^{}n)\ −\ r\ U\^{}n},
  where \(T2 = tridiag(1,−2,1)\) is the second-difference and
  \(T1 = tridiag(−1,0,1)\) the central first-difference. Solving for
  \(U^{n+1}\) gives
  \texttt{U\^{}\{n+1\}\ =\ {[}(1−rk)I\ +\ ½kσ²/h²\ T2\ +\ k(r−½σ²)/(2h)\ T1{]}\ U\^{}n\ =\ F\ U\^{}n}.
  Identification: \((1−rk)I\) is the \textbf{reaction} (the \(−rU\)
  discount term), \texttt{½kσ²/h²\ T2} is the \textbf{diffusion}, and
  \texttt{k(r−½σ²)/(2h)\ T1} is the \textbf{advection} (drift of
  log-price).
\item
  \textbf{\texttt{:29-33}} --- \(L = 2 log K + 2\), \(Nx = 1000\),
  \(Nt = 40000*T\). Why center the grid at \(log K\) (\texttt{:43})?
  What is the analogy between this domain truncation and the COS
  \([a,b]\) range bug you hit in Assignment 2? \textbf{A.} The payoff
  kink and the early-exercise boundary both sit near \(S=K\),
  i.e.~\(m=ln K\), so centring the grid there
  (\texttt{mvec\ +=\ np.log(K)}, line 43) puts the finest action and the
  interpolation point \(ln S0\) in the well-resolved interior, keeping
  the truncated boundaries far from the region of interest. The analogy:
  in both methods you replace an integral/PDE on \((−∞,∞)\) by one on a
  finite \([a,b]\); if \([a,b]\) (here \(L\)) is too narrow it cuts off
  mass/payoff support and the boundary error dominates and \emph{will
  not} vanish as you refine \(N\) (COS) or \(Nx,Nt\) (FD). In Assignment
  2 the COS \([a,b]\) collapsed at short maturity (small \(c2=σ²T\));
  here \(L = 2ln K + 2\) is fixed and generous, the FD analogue of
  choosing \(L√(c2+√c4)\) wide enough. \emph{(confirm this matches your
  own reasoning)}
\item
  \textbf{\texttt{:45-46} and \texttt{:51}} --- initial condition
  \(max(K - e^m, 0)\) and the boundary term \(p[0]\) scaled by
  \(K·exp(-r·time2mat)\). What Dirichlet boundary condition does
  \(p[0]\) enforce at the deep-ITM edge, and why is no correction added
  at the high-\(S\) boundary? \textbf{A.} The initial condition (line
  46, \(τ=0\)) is the put payoff \(max(K−e^m,0)\). At the lower edge
  \(m→−∞\) (deep ITM, \(S→0\)) the put value tends to
  \texttt{K\ e\^{}\{−rτ\}\ −\ S\ ≈\ K\ e\^{}\{−rτ\}} (discounted
  strike). The central-difference stencil at the first interior node
  references the ghost node just below the grid, whose Dirichlet value
  is \(K e^{−rτ}\); \(p[0]\) (line 51) injects exactly that boundary
  value, scaled by the matching stencil coefficient
  \texttt{(½kσ²/h²\ −\ k(r−½σ²)/(2h))} and \(K·exp(−r·time2mat)\). At
  the high-\(S\) edge (\(m→+∞\)) the put value → 0, so the Dirichlet
  value there is zero and no correction term is needed --- \(p\) stays
  zero on that side.
\item
  \textbf{\texttt{:105}} --- \texttt{U\ =\ np.maximum(U,\ intrinsic)},
  ``the single line that makes it American''. Why is projecting onto the
  payoff after each explicit step a valid (operator-splitting /
  explicit-penalty) treatment of the free boundary, and what accuracy
  does it cost versus a proper LCP/PSOR solve? \textbf{A.} The American
  put solves a linear complementarity problem: \(U ≥ g\), \(LU ≤ 0\),
  with complementarity. The explicit projection \(U ← max(U, g)\) after
  each time step is a Bermudan-style operator splitting: advance the PDE
  one step, then enforce the constraint \(U ≥ intrinsic\) by projecting
  the violated nodes up to the payoff. Because the time step \(k\) is
  tiny, the splitting error per step is \(O(k)\) and the discrete
  exercise dates approximate continuous exercise. It is valid but only
  \emph{first-order} accurate in the constraint handling and it smears
  the free boundary by one grid/time cell. A proper LCP/PSOR solve
  treats the constraint \emph{implicitly and simultaneously} with the
  PDE at each step (projected SOR iteration), locating the free boundary
  more accurately and without the explicit-projection bias --- at the
  cost of an iterative solve per step. \emph{(confirm this matches your
  own reasoning)}
\item
  \textbf{Stability:} with \(Nx=1000\), \(Nt=40000·T\), worst-case
  \(σ=0.4\), compute the diffusion number \texttt{0.5\ σ²\ k\ /\ h²} and
  check it is \texttt{≤\ 1/2}. Is the scheme stable? Why might the
  author have picked \(Nt\) so large (stability vs accuracy)? Can the
  \(T1\) advection term cause oscillations even when the diffusion bound
  holds? \textbf{A.} \texttt{h\ =\ L/Nx} with \(L = 2ln40+2 ≈ 9.378\),
  so \texttt{h\ ≈\ 9.378/1000\ ≈\ 9.38e-3}, \(h² ≈ 8.79e-5\).
  \texttt{k\ =\ T/Nt\ =\ T/(40000T)\ =\ 1/40000\ =\ 2.5e-5}. Then
  \texttt{½σ²k/h²\ =\ ½·0.16·2.5e-5/8.79e-5\ ≈\ 0.5·0.16·0.2845\ ≈\ 0.0228},
  well under \(½\). So the explicit FTCS scheme is stable (von Neumann
  diffusion bound satisfied with large margin). \(Nt\) is huge mostly to
  satisfy this \texttt{k\ ≤\ h²/σ²} stability constraint --- explicit
  schemes force \(k = O(h²)\), so refining \(Nx\) (smaller \(h\))
  demands quadratically more time steps; accuracy is a secondary
  beneficiary. The \(T1\) advection (central difference) can produce
  spurious oscillations when the \emph{cell Péclet number}
  \texttt{\textbar{}r−½σ²\textbar{}h/σ²} exceeds \textasciitilde1,
  independently of the diffusion bound; here drift is small and \(h\)
  tiny so Péclet \textless\textless{} 1, no oscillation, but in
  principle yes it can. \emph{(confirm the numbers against your
  printout)}
\item
  \textbf{\texttt{:54} / \texttt{:107}} --- \texttt{np.interp} reads the
  price at \(log(S0)\). Your FD-European prices are systematically
  \(~1e-4\) \emph{below} Black--Scholes (every difference negative). Is
  the consistent sign due to the explicit scheme, the truncated domain,
  or the interpolation? Defend your pick. \textbf{A.} The consistent
  negative sign points to the \emph{truncated domain} as the dominant
  source. Linear \texttt{np.interp} on a convex value function biases
  the read-off \emph{upward} (chord lies above the curve), so
  interpolation alone would push FD \emph{above} BS --- wrong sign. The
  explicit FTCS truncation error is \(O(k)+O(h²)\) and is essentially
  unsigned/oscillatory across the parameter grid, so it would not
  produce a uniform sign on every row. Truncating the domain removes a
  sliver of probability mass in the tails that the option payoff would
  have weighted, so the discrete value is systematically a touch
  \emph{low} everywhere --- a one-sided, consistent \textasciitilde1e-4
  deficit, matching the observation. So I attribute the uniform negative
  difference to domain truncation, not to the scheme or the
  interpolation. \emph{(confirm this matches your own observation of the
  sign on every row)}
\end{enumerate}

\hypertarget{assignment-set-2}{%
\section{Assignment set 2}\label{assignment-set-2}}

Files: - \texttt{assignments/assignment-set-2/Q1\_basket\_MC\ (1).py} -
\texttt{assignments/assignment-set-2/Q2\_moment\_matching\ (1).py} -
\texttt{assignments/assignment-set-2/Q3\_basket\_COS\ (1).py}

\hypertarget{q3-basket-cos-highest-priority-fang-invented-cos}{%
\subsection{Q3 --- basket COS (highest-priority: Fang invented
COS)}\label{q3-basket-cos-highest-priority-fang-invented-cos}}

\begin{enumerate}
\def\labelenumi{\arabic{enumi}.}
\item
  \textbf{\texttt{Q3\_basket\_COS\ (1).py:22-32}} --- \(chi_coeffs\):
  \texttt{chi\_k\ =\ Re\{\ phi(omega\_k)\ exp(-i\ omega\_k\ a)\ \}},
  \texttt{omega\_k\ =\ k\ pi/(b-a)}, with \(chi[0] *= 0.5\) at
  \texttt{:31}. Explain term by term where this comes from, and why the
  \(k=0\) term is halved (the primed sum). \textbf{A.} These are the
  \emph{cosine} density coefficients of the COS method. Expanding the
  density on \([a,b]\) in \texttt{\{cos(kπ(x−a)/(b−a))\}}, the
  coefficient is
  \texttt{A\_k\ =\ (2/(b−a))\ ∫\_a\^{}b\ f(x)\ cos(ω\_k(x−a))\ dx},
  \texttt{ω\_k\ =\ kπ/(b−a)} (line 29). Writing
  \texttt{cos(ω\_k(x−a))\ =\ Re\{e\^{}\{iω\_k(x−a)\}\}} and using
  \texttt{∫\ f(x)e\^{}\{iω\_k\ x\}dx\ =\ φ(ω\_k)} gives
  \texttt{A\_k\ ∝\ Re\{φ(ω\_k)\ e\^{}\{−iω\_k\ a\}\}} --- exactly line
  30. (The code folds the \texttt{2/(b−a)} and the payoff normalisation
  into the final triple product rather than here.) The \(k=0\) term is
  halved (line 31, \(chi[0] *= 0.5\)) because the COS reconstruction
  \(f(x) ≈ Σ' A_k cos(...)\) uses a \emph{primed} sum where the first
  term carries weight ½ --- this is the standard Fourier-cosine
  convention (the constant mode is the average, counted once, not
  twice).
\item
  \textbf{\texttt{:34-54}} --- \texttt{payoff\_coeffs\_2d} builds
  \(V_{k1,k2}\) as \(C1^T P C2\) by trapezoidal quadrature on a dense
  payoff grid. Why is there \emph{no} closed-form \(V_k\) here, unlike
  the 1D European call/put (where \texttt{χ\_k,\ ψ\_k} are analytic)?
  What couples the two state variables? \textbf{A.} In 1D the payoff
  \(max(±(S−K),0)\) is a function of a \emph{single} variable and the
  cosine integral of \(e^y\) and \(1\) against \(cos\) has the
  closed-form \texttt{χ\_k,\ ψ\_k}. Here the payoff is
  \(max(w1 g1(y1) + w2 g2(y2) − K, 0)\) --- the \(max(·,0)\) and the
  strike couple \(y1\) and \(y2\) \emph{non-separably}: the region where
  the basket exceeds \(K\) is a diagonal half-plane
  \({w1 g1(y1)+w2 g2(y2) > K}\), not a product of intervals, so the 2D
  integral does not factor into a product of 1D integrals and has no
  elementary antiderivative. Hence line 49 builds the payoff on a dense
  grid \(P\) and lines 52--54 apply 2D trapezoidal quadrature written as
  \(C1ᵀ P C2\), where \(C1, C2\) are the cosine matrices times the
  trapezoidal weights \((Wy·h)\). It is the strike \(K\) inside the
  joint \(max\) that couples the two states.
\item
  \textbf{\texttt{:56-63}} --- \(price_2dCOS\) assembles
  \texttt{e\^{}\{-rT\}\ ·\ chi1\^{}T\ V\ chi2}. What makes the 2D
  integral collapse to this triple product of coefficient vectors, and
  where does the \(rho=0\) independence assumption enter (the density
  coefficients factorising
  \texttt{c\_\{k1,k2\}=chi\_\{k1\}\ chi\_\{k2\}})? \textbf{A.} The price
  is \(e^{−rT} ∫∫ payoff(y1,y2) f(y1,y2) dy1 dy2\). Expanding the
  \emph{joint} density in a 2D cosine series,
  \texttt{f(y1,y2)\ ≈\ Σ\_\{k1,k2\}\ c\_\{k1,k2\}\ cos(ω\_\{k1\}(y1−a1))\ cos(ω\_\{k2\}(y2−a2))},
  and by 2D Parseval the integral becomes
  \texttt{Σ\_\{k1,k2\}\ c\_\{k1,k2\}\ V\_\{k1,k2\}\ =\ chi1ᵀ\ V\ chi2}
  once \(c_{k1,k2}\) factorises. The \(rho=0\) independence is exactly
  what makes it factorise: with independent marginals the joint density
  is a product \(f(y1,y2)=f1(y1)f2(y2)\), so its 2D cosine coefficients
  separate \texttt{c\_\{k1,k2\}\ =\ chi\_\{k1\}\ chi\_\{k2\}} (the outer
  product of the two 1D coefficient vectors from \(chi_coeffs\)). That
  is why the triple product \(chi1ᵀ V chi2\) --- and not a full
  coefficient \emph{matrix} recovered from a 2D characteristic function
  --- suffices. (If ρ≠0 you would need the joint CF and a 2D coefficient
  matrix.)
\item
  \textbf{\texttt{:65-71}} --- \(cumulant_range\) returns
  \([k1 - width, k1 + width]\) with
  \(width = L*sqrt(|k2| + sqrt(|k4|))\), \(L=12\). Interpret each of
  \(k1\), \(k2\), \(sqrt(k4)\), \(L\). \textbf{A.} \(k1\) is the first
  cumulant = mean of the marginal, so the interval is centred on the
  mean. \(k2\) is the second cumulant = variance, so \(√k2\) is the
  standard deviation --- the basic width scale. \(√k4\) adds the
  fourth-cumulant (kurtosis/tail) contribution: \(k4>0\) (heavy tails)
  widens the interval to capture the extra tail mass that a
  pure-Gaussian \(√k2\) would miss. \(L\) is the number of
  ``standard-deviation-like units'' of half-width (here \(L=12\)), the
  safety multiplier from Fang--Oosterlee; larger \(L\) means a wider,
  safer truncation at the cost of needing more terms \(N\) for the same
  resolution.
\item
  \textbf{\texttt{:82-83}} --- Q3(a) calls
  \(cumulant_range(mu, sigma^2 T, 0)\) per marginal. For a normal
  log-price, what are \(k1\), \(k2\), and why is \(k4 = 0\)? \textbf{A.}
  For the normal log-price
  \texttt{x\_i\ =\ ln\ S\_i(T)\ \textasciitilde{}\ N(μ\_i,\ σ\_i²T)}
  with \texttt{μ\_i\ =\ ln100\ +\ (r−½σ\_i²)T} (lines 76--77), the first
  cumulant \(k1 = μ_i\) (the mean) and the second cumulant
  \(k2 = σ_i²T\) (the variance). A Gaussian has \emph{every} cumulant of
  order ≥3 equal to zero, so \(k4 = 0\) exactly --- there is no excess
  kurtosis, the tails are exactly Gaussian, and the width reduces to
  \texttt{L\ σ\_i√T\ =\ 12σ\_i√T}.
\item
  \textbf{THE BUG (over-prepare this).} \texttt{:70} ---
  \(width = L*sqrt(|k2| + sqrt(|k4|))\). As \(T→0\),
  \(k2 = sigma^2 T → 0\), so \(width\) collapses and \([a,b]\) becomes
  too narrow to contain the density/payoff support. Show me the exact
  line where this happens. \textbf{TODO (Bruno):} what symptom did you
  see and at which \(T\); and exactly what did you change --- \(L\) at
  \texttt{:65}, the width formula at \texttt{:70}, or the non-negativity
  clamp at \texttt{:127}? Why is \(sqrt(c2 + sqrt(c4))\) (not just
  \(sqrt(c2)\)) the right safeguard, and why is a fixed wide interval
  not a free lunch? \textbf{A.} The collapse is on \textbf{line 70}:
  \texttt{width\ =\ L*np.sqrt(abs(k2)\ +\ np.sqrt(abs(k4)))}. For the
  normal legs \(k2 = σ²T\) and \(k4 = 0\), so \(width = Lσ√T\), which
  \(→ 0\) as \(T→0\); with \(k4=0\) there is \emph{no floor} on the
  width, so \([a,b] = [μ − Lσ√T, μ + Lσ√T]\) shrinks to a point around
  the mean and no longer brackets the support of the density or,
  crucially, the payoff kink at the strike. The mechanism: at short
  \(T\) the truncation interval can become narrower than the distance
  from the mean to \(ln K\), so the diagonal exercise region
  \({basket > K}\) is partly outside \([a,b]\) and the quadrature on the
  dense grid (lines 43--49) integrates a payoff grid that misses the
  in-the-money tail. What one \emph{would} expect to see as the symptom:
  the COS price stops matching the MC/moment benchmark at small
  maturities (e.g.~\(T\) of order 0.05--0.1), the error \emph{fails to
  fall} as you raise \(N\) (because it is a domain-truncation error, not
  a series error), and the price can even drift or go visibly wrong
  while the long-maturity prices stay correct. The standard fixes, in
  order of preference: keep the \(√(c2 + √c4)\) form (do \textbf{not}
  drop to \(√c2\)) and/or add an explicit floor on the half-width, or
  raise \(L\) (line 65); the cleanest is a width floor so \([a,b]\)
  always contains \(ln K ± a few σ\). The \(√(c2 + √c4)\) form is the
  right safeguard because for heavy-tailed/short-maturity laws the
  variance \(c2\) alone \emph{understates} how far the relevant mass
  reaches --- the \(√c4\) term adds the kurtosis/tail contribution and
  keeps the interval from collapsing when \(c2\) is tiny but the tails
  (or the strike distance) still matter. A fixed wide interval is not a
  free lunch because COS resolution is \texttt{N/(b−a)}: widening
  \([a,b]\) for a fixed \(N\) \emph{coarsens} the cosine grid, so you
  trade truncation error for series-truncation error and must raise
  \(N\) (more cost) to keep the same accuracy --- the cumulant rule is
  precisely the device that picks \([a,b]\) just wide enough.
  \textbf{(Bruno: confirm which exact change you made --- \(L\) at line
  65, the width formula at line 70, or a floor --- and at which \(T\)
  you saw the price break; the analysis above is the correct general
  justification but the specific number/edit is yours.)}
\item
  \textbf{\texttt{:117-119}} --- \(phi2_b\) is the CIR-leg
  characteristic function \texttt{exp(lambda·i·z/(1-2·i·z))},
  \(z = om·c\). Where does this come from (the zero-dof non-central
  chi-square law of \(S2(T)=c·Z\))? \textbf{A.} By the technical hint,
  the CIR terminal value factors as \(S2(T) = c·Z\) with
  \(Z ~ χ'²_0(λ)\), a non-central chi-square with \textbf{zero} degrees
  of freedom and non-centrality \(λ\) (lines 114--116 set \(c_b\),
  \(lam_b\)). The CF of a non-central chi-square \(χ'²_ν(λ)\) is
  \texttt{exp(iλt/(1−2it))\ /\ (1−2it)\^{}\{ν/2\}}; with \(ν=0\) the
  denominator power vanishes, leaving
  \texttt{φ\_Z(t)\ =\ exp(iλt/(1−2it))}. Since \(S2(T) = cZ\),
  \texttt{φ\_\{S2\}(ω)\ =\ φ\_Z(cω)}, i.e.~with \(z = ω·c\) the CF is
  \texttt{exp(λ·iz/(1−2iz))} --- exactly line 119. The zero-dof part (no
  \texttt{(1−2it)\^{}\{ν/2\}} factor) is the special structure of the
  square-root process at this parameterisation.
\item
  \textbf{\texttt{:122-127}} --- CIR cumulants \texttt{k1\ =\ c·lambda},
  \texttt{k2\ =\ 4\ c\^{}2\ lambda},
  \texttt{k4\ =\ 192\ c\^{}4\ lambda}, and the range is clamped
  \(a2 = max(a2, 0)\) at \texttt{:127}. Why is the clamp needed for the
  square-root process? \textbf{A.} For \(Z ~ χ'²_0(λ)\) the cumulants
  are \(κ_n = 2^{n−1} n! λ\) (so \(κ1=λ\), \(κ2=4λ\), \(κ4=192λ\)), and
  since \(S2=cZ\) they scale as \(c^n\): \(k1=cλ\), \(k2=4c²λ\),
  \(k4=192c⁴λ\) (lines 122--124). The cumulant rule
  \([k1 ± L√(k2+√k4)]\) is \emph{symmetric} about the mean, but
  \(S2(T)\) is a CIR/square-root level that is \textbf{non-negative} by
  construction --- the process cannot go below 0. For the heavily
  right-skewed CIR law the symmetric lower endpoint \(a2 = k1 − width\)
  can fall below 0, which is outside the support and would waste grid
  (and put the quadrature on a region of zero density). The clamp
  \(a2 = max(a2, 0)\) (line 127) restricts the truncation interval to
  the physical support \([0, b2]\), matching the non-negativity of the
  square-root process.
\end{enumerate}

\hypertarget{q1-basket-monte-carlo}{%
\subsection{Q1 --- basket Monte Carlo}\label{q1-basket-monte-carlo}}

\begin{enumerate}
\def\labelenumi{\arabic{enumi}.}
\setcounter{enumi}{8}
\item
  \textbf{\texttt{Q1\_basket\_MC\ (1).py:29} and \texttt{:116}} ---
  \(Z2 = rho*Z1 + sqrt(1-rho^2)*Zc\). Show this is exactly the 2×2
  Cholesky factor of the correlation matrix, i.e.~that
  \(Corr(Z1,Z2)=rho\). \textbf{A.} \(Z1, Zc ~ N(0,1)\) independent.
  \(Var(Z2) = ρ²Var(Z1) + (1−ρ²)Var(Zc) = ρ² + (1−ρ²) = 1\), and
  \(Cov(Z1,Z2) = E[Z1(ρZ1 + √(1−ρ²)Zc)] = ρE[Z1²] + √(1−ρ²)E[Z1 Zc] = ρ·1 + 0 = ρ\).
  So \texttt{Corr(Z1,Z2)\ =\ ρ/(1·1)\ =\ ρ}. The transform
  \((Z1,Zc) ↦ (Z1,Z2)\) is the lower-triangular matrix
  \(L = [[1,0],[ρ,√(1−ρ²)]]\), which is precisely the Cholesky factor of
  \(Σ = [[1,ρ],[ρ,1]] = L Lᵀ\).
\item
  \textbf{\texttt{:120-121}} --- Euler step for the CIR leg with full
  truncation \(S2p = max(S2,0)\) inside both drift and diffusion.
  \textbf{TODO (Bruno):} why full truncation rather than reflection, and
  what bias does it introduce? Why is the Q-drift \(r·S2\) (not the
  physical \(kappa(theta-S2)\))? \textbf{A.} Lines 120--121 apply
  \(S2p = max(S2,0)\) and step \(S2 = S2 + r·S2p·h + σ2·√S2p·√h·Z2\).
  The Euler discretisation of a square-root diffusion can overshoot
  below zero in a step, leaving \(√S2\) undefined; \textbf{full
  truncation} (Lord, Koekkoek \& van Dijk 2010) substitutes \(S2⁺\) only
  \emph{inside the coefficients} (drift and diffusion) while letting the
  state itself go slightly negative, then truncates. It is preferred
  over \textbf{reflection} (\(|S2|\) or \(−S2\)) because full truncation
  is the scheme with the smallest discretisation bias among the fixes in
  that paper and it does not artificially inject spurious positive
  drift/variance that reflection creates by bouncing paths off zero;
  reflection systematically biases the level upward near the boundary.
  Full truncation still introduces a (downward) bias because zeroing the
  coefficient kills the local drift/vol whenever a path is at/below
  zero, but it is the mildest. The Q-drift is \(r·S2\), not the physical
  mean-reverting \(κ(θ−S2)\), because under the risk-neutral measure
  \(S2\) is a \emph{traded asset} whose discounted price \(e^{−rt}S2\)
  must be a martingale; that forces drift \(r·S2 dt\). The physical CIR
  drift \(κ(θ−S2)\) belongs to a (non-traded) rate/variance factor under
  the real-world measure --- here \(S2\) is a price, so risk-neutral
  valuation pins its drift to \(r·S2\). \emph{(confirm this matches your
  own reasoning on the full-truncation choice and any bias you
  observed)}
\item
  \textbf{\texttt{:144-169}} --- the weak-error test reuses the
  \emph{same} fine Brownian increments across all coarse grids. Why does
  reusing increments cleanly isolate the discretization bias from the
  Monte-Carlo noise? \textbf{A.} The total error of a discretised MC
  price is \texttt{(discretisation\ bias)\ +\ (statistical/MC\ noise)}.
  By generating one set of fine Brownian increments \(dW1f, dW2f\) per
  chunk and \emph{aggregating} them into each coarse grid
  (\texttt{basket\_on\_grid} sums \texttt{k\ =\ n\_fine/n\_steps} fine
  increments into each coarse step, lines 149--152), every grid is
  driven by the \emph{same} underlying Brownian path. So when you
  difference two grids, the MC noise (which depends only on the realised
  path) cancels in common, and what remains is purely the difference in
  \emph{discretisation} --- the bias due to step size \(h\). The price
  difference \(p(n) − p_ref\) is then an almost noise-free estimate of
  the weak error, which is why the fitted slope on lines 181--185
  cleanly recovers the expected weak order 1.0 for Euler; without common
  increments the MC noise (\textasciitilde{}\texttt{1/√M}) would swamp
  the small \(O(h)\) bias.
\end{enumerate}

\hypertarget{q2-moment-matching}{%
\subsection{Q2 --- moment matching}\label{q2-moment-matching}}

\begin{enumerate}
\def\labelenumi{\arabic{enumi}.}
\setcounter{enumi}{11}
\item
  \textbf{\texttt{Q2\_moment\_matching\ (1).py:20-29}} ---
  \(lognormal_cross\) returns \texttt{E{[}S1\^{}j\ S2\^{}m{]}} via
  \(exp(mean + 0.5 var)\). Derive this cross raw moment for two
  correlated log-normals by hand. \textbf{A.} Write \(S1 = e^X\),
  \(S2 = e^Y\) with \((X,Y)\) jointly normal:
  \texttt{X\ \textasciitilde{}\ N(μ\_X,\ v\_X)},
  \texttt{Y\ \textasciitilde{}\ N(μ\_Y,\ v\_Y)}, \(Cov(X,Y) = c_{XY}\)
  (lines 26--28: \(μ_X = ln S10 + (r−½σ1²)T\), \(v_X = σ1²T\),
  \(c_{XY} = ρσ1σ2T\)). Then
  \texttt{S1\^{}j\ S2\^{}m\ =\ e\^{}\{jX\ +\ mY\}}, and \(jX + mY\) is
  normal with mean \texttt{jμ\_X\ +\ mμ\_Y} and variance
  \texttt{j²v\_X\ +\ m²v\_Y\ +\ 2jm\ c\_\{XY\}}. Using the MGF of a
  normal, \(E[e^N] = exp(mean + ½var)\) for \(N\) normal, gives
  \texttt{E{[}S1\^{}j\ S2\^{}m{]}\ =\ exp(jμ\_X\ +\ mμ\_Y\ +\ ½(j²v\_X\ +\ m²v\_Y\ +\ 2jm\ c\_\{XY\}))}
  --- exactly line 29.
\item
  \textbf{\texttt{:39-53}} --- \texttt{cir\_S2\_rawmoments} uses
  cumulants \texttt{kappa\_n\ =\ 2\^{}\{n-1\}\ n!\ lambda} then a
  raw-from-cumulant recursion. Derive the non-central chi-square
  (zero-dof) cumulants \(kappa_n\). \textbf{A.} From
  \texttt{φ\_Z(t)\ =\ exp(λ·it/(1−2it))} (zero-dof non-central
  chi-square, see Q3 Q7), the cumulant generating function is
  \texttt{K(t)\ =\ log\ E{[}e\^{}\{tZ\}{]}\ =\ λt/(1−2t)} (set
  \(it → t\)). Expand \texttt{1/(1−2t)\ =\ Σ\_\{m≥0\}(2t)\^{}m}, so
  \texttt{K(t)\ =\ λ\ Σ\_\{m≥0\}\ 2\^{}m\ t\^{}\{m+1\}\ =\ λ\ Σ\_\{n≥1\}\ 2\^{}\{n−1\}\ t\^{}n}.
  The cumulant \(κ_n\) is \(n!\) times the coefficient of \(t^n\) in
  \(K(t)\), hence
  \texttt{κ\_n\ =\ n!·λ·2\^{}\{n−1\}\ =\ 2\^{}\{n−1\}\ n!\ λ} (line 49).
  The code then turns these cumulants into raw moments by the standard
  recursion
  \texttt{m\_n\ =\ Σ\_\{k=0\}\^{}\{n−1\}\ C(n−1,k)\ κ\_\{n−k\}\ m\_k}
  (lines 51--52), and scales \(S2 = cZ\) by \(c^{k+1}\) (line 53).
\item
  \textbf{\texttt{:78-139}} --- \(price_2moments\) (log-normal),
  \(price_3moments\) (shifted log-normal via skewness inversion with
  \(brentq\)), \(price_4moments\) (Johnson SU via \(least_squares\)).
  \textbf{TODO (Bruno):} why is each family the natural choice for 2/3/4
  matched moments, and why need not matching more moments monotonically
  reduce the pricing error? \textbf{A.} Each family has exactly as many
  free shape parameters as moments to match, and is the canonical
  positive-support distribution at that order: (2) a \textbf{log-normal}
  has two parameters \((μ,σ²)\) → matches mean and variance, and is
  natural because the basket is a positive sum close to log-normal
  (Black/Lévy moment-matching); skewness/kurtosis are then
  \emph{implied}, not controlled. (3) a \textbf{shifted log-normal}
  \(γ + L\) adds a location shift \(γ\), giving three parameters, so it
  can also match \textbf{skewness} --- the code inverts the closed-form
  log-normal skewness \((w+2)√(w−1) = γ1\) for \(w = e^{s²}\) with
  \(brentq\) (line 99), then sets \(M, γ\) from variance and mean. (4)
  the \textbf{Johnson SU} \texttt{ξ\ +\ λ\ sinh((Z−γ)/δ)} has four
  parameters \((ξ,λ,δ,γ)\), enough to match mean, variance,
  \textbf{skewness and kurtosis}; its \((δ,γ)\) are fitted to the
  standardized skew/kurtosis by \(least_squares\) (line 134) because
  their closed forms are awkward, then \((ξ,λ)\) from mean/variance.
  Matching more moments need \emph{not} reduce the pricing error
  monotonically because: (a) the call price is an integral against the
  \emph{whole} density tail, and the error depends on how well the
  chosen family's shape matches the true basket density near and beyond
  the strike --- a higher-moment family can fit the body better yet
  distort the tail that the payoff weights; (b) the moment problem can
  be ill-conditioned (e.g.~the Johnson \(least_squares\) may land on a
  shape whose tail differs from the true one even though skew/kurtosis
  agree); and (c) the basket's true law is not exactly any of these
  families, so adding a constraint trades one mismatch for another
  rather than strictly improving. So 3- or 4-moment can occasionally do
  \emph{worse} than 2-moment on the price even though it matches more
  moments. \emph{(confirm this matches the non-monotone error pattern
  you saw in your Q2 table)}
\end{enumerate}

\end{document}
